Matematične dokaze iz logike lahko interpretiramo kot računalniške programe ali pa kot sklepe v teoriji tipov. 
Notacija je zelo podobna, kot ta ki smo jo obravnavali v prejšnjem poglavju. Spodaj je primerjava pravil sklepanja.


\begin{align*}
    \begin{bprooftree}
        \AXC{\(\G \vdash A\) \tru}
        \AXC{\(\G \vdash B\) \tru}
        \RL{\(\land I\).} %\mathcal{I}
        \BIC{\(\G \vdash A \land B\) \tru}
    \end{bprooftree}
    \qquad
    \begin{bprooftree}
        \AXC{\(\G \vdash a:A\)}
        \AXC{\(\G \vdash b:B\)}
        \RL{\(\times I\).}
        \BIC{\(\G \vdash (a,b): A \times B\)}
    \end{bprooftree}
\end{align*}

\begin{align*}
    \begin{bprooftree}
        \AXC{\(A \land B\) \tru}
        \RL{\(\land E_L\)}
        \UIC{\(A\) \tru}
    \end{bprooftree}
    \qquad&
     \begin{bprooftree}
        \AXC{\(A \land B\) \tru}
        \RL{\(\land E_D\)}
        \UIC{\(A\) \tru}
    \end{bprooftree}
    \qquad&
    \begin{bprooftree}
        \AXC{\(x: A \times B\)}
        \RL{\(\times E_1\)}
        \UIC{\(fst(x):A\)}
    \end{bprooftree}
    \qquad&
     \begin{bprooftree}
        \AXC{\(x: A \times B\)}
        \RL{\(\times E_2\)}
        \UIC{\(snd(x) : A\)}
    \end{bprooftree}
\end{align*}
\begin{align*}
    \begin{bprooftree}
        \AXC{\(\)}
        \RL{\(u\)}
        \UIC{\(A\) \tru}
        \noLine
        \UIC{\(\vdots\)}
        \noLine
        \UIC{\(B\) \tru}
        \RL{\(\Rightarrow I^u\)}
        \UIC{\(A \implies B\) \tru.}
    \end{bprooftree}
    \qquad&
    \begin{bprooftree}
        \AXC{\(\)}
        \RL{\(u\)}
        \UIC{\(a: A\)}
        \noLine
        \UIC{\(\vdots\)}
        \noLine
        \UIC{\(b: B\)}
        \RL{\(\rightarrow I^u\)}
        \UIC{\(f: A \rightarrow B\).}
    \end{bprooftree}    
    \\\\
    \begin{bprooftree}
        \AXC{\(A \implies B\) \tru}
        \AXC{\(A\) \tru}
        \RL{\(\Rightarrow E\)}
        \BIC{\(B\) \tru}
    \end{bprooftree}
    \qquad&
    \begin{bprooftree}
        \AXC{\(f: A \rightarrow B\)}
        \AXC{\(a: A\)}
        \RL{\(\rightarrow E\)}
        \BIC{\(f(a) : B\)}
\end{bprooftree}
\end{align*}

% Disjunkcija
\begin{align*}
    \begin{bprooftree}
        \AXC{\(A\) \tru}
        \RL{\(\lor I_L\)}
        \UIC{\(A \lor B\) \tru}
    \end{bprooftree}
    \qquad&
    \begin{bprooftree}
        \AXC{\(B\) \tru}
        \RL{\(\lor I_D\)}
        \UIC{\(A \lor B\) \tru}
    \end{bprooftree}
    \qquad&
    %
    % TODO: vsotni tip
    \begin{bprooftree}
        \AXC{\(a: A\)}
        \RL{\(\lor I_L\)}
        \UIC{\(A \lor B\)}
    \end{bprooftree}
    \qquad&
    \begin{bprooftree}
        \AXC{\(B\)}
        \RL{\(\lor I_D\)}
        \UIC{\(A \lor B\)}
    \end{bprooftree}
\end{align*}

\begin{align*}
    \begin{bprooftree}
    \AXC{\(A \lor B\)}
%
    \AXC{}
    \RL{\(u\)}
    \UIC{\(A\) \tru}
    \noLine
    \UIC{\(\vdots\)}
    \noLine
    \UIC{\(C\) \tru}
%
    \AXC{}
    \RL{\(w\)}
    \UIC{\(B\) \tru}
    \noLine
    \UIC{\(\vdots\)}
    \noLine
    \UIC{\(C\) \tru}
    \RL{\(\lor E^{u, w}\) }
    \TIC{\(C\) \tru}
\end{bprooftree}
\qquad&
\begin{bprooftree}
    \AXC{\(A \lor B\)}
%
    \AXC{}
    \RL{\(u\)}
    \UIC{\(A\)}
    \noLine
    \UIC{\(\vdots\)}
    \noLine
    \UIC{\(C\)}
%
    \AXC{}
    \RL{\(w\)}
    \UIC{\(B\)}
    \noLine
    \UIC{\(\vdots\)}
    \noLine
    \UIC{\(C\)}
    \RL{\(\lor E^{u, w}\) }
    \TIC{\(C\)}
\end{bprooftree}
\end{align*}